\chapter{Getaltheorie}\label{ch:g12:arith}

De getaltheorie bestudeert de \omterm{def:g10:numbers:sets}{gehele getallen}: \omterm{def:g12:arith:divides}{deelbaarheid}, priemgetallen,
resten. Lang gold ze als de zuiverste van de zuivere wiskunde; vandaag
beschermt ze elke betaling op het internet: het RSA-systeem steunt op de
stellingen van Bézout, Gauss en Fermat die in dit hoofdstuk bewezen worden.

\section{Deelbaarheid en euclidische deling}

\begin{definition}[Deelbaarheid]\label{def:g12:arith:divides}
Zij $a, b \in \Z$. We zeggen dat $b$ het getal $a$ \emph{deelt}\index{deelbaarheid},
genoteerd $b \mid a$, als er een $k \in \Z$ bestaat met $a = kb$. We zeggen
dan ook dat $a$ een \emph{veelvoud} van $b$ is.
\end{definition}

\begin{proposition}\label{prop:g12:arith:divprops}
Geldt $c \mid a$ en $c \mid b$, dan \omterm{def:g12:arith:divides}{deelt} $c$ elke gehele combinatie
$au + bv$ ($u, v \in \Z$). Geldt $a \mid b$ en $b \mid a$ met
$a,b \in \N$, dan is $a = b$. Geldt $a \mid b$ en $b \neq 0$, dan is
$\abs a \leq \abs b$.
\end{proposition}

\begin{proof}
Schrijf $a = kc$ en $b = lc$: dan is $au + bv = (ku + lv)c$. De overige punten
volgen uit $\abs{a} = \abs{k}\,\abs{b}$ met $\abs k \geq 1$ zodra
$b = ka \neq 0$.
\end{proof}

\begin{theorem}[Euclidische deling]\label{thm:g12:arith:euclid}
\index{euclidische deling}
Zij $a \in \Z$ en $b \in \N^*$. Er bestaat precies één paar
$(q, r) \in \Z \times \N$ met
\[
a = bq + r \qquad\text{en}\qquad 0 \leq r < b .
\]
Hierbij is $q$ het \emph{quotiënt} en $r$ de \emph{rest}.
\end{theorem}

\begin{proof}
\emph{Bestaan.} De verzameling van de veelvouden van $b$ die $a$ niet
overschrijden heeft een grootste element $bq$ (ze is niet leeg en \omterm{def:g12:seq:bounded}{naar boven
begrensd}); stel $r = a - bq$. Wegens de maximaliteit is $b(q+1) > a$, dus
$0 \leq r < b$.
\emph{Uniciteit.} Geldt $bq + r = bq' + r'$ met $0 \leq r, r' < b$, dan is
$b(q - q') = r' - r$ met $\abs{r' - r} < b$: een veelvoud van $b$ met \omterm{def:g10:numbers:abs}{absolute
waarde} kleiner dan $b$ is noodzakelijk $0$, dus $r = r'$ en $q = q'$.
\end{proof}

\section{Congruenties}

\begin{definition}[Congruentie]\label{def:g12:arith:congruence}
Zij $n \in \N^*$. Twee \omterm{def:g10:numbers:sets}{gehele getallen} $a$ en $b$ heten \emph{congruent modulo
$n$}\index{congruentie}, genoteerd $a \equiv b \pmod n$, als $n \mid (a - b)$
--- gelijkwaardig: als $a$ en $b$ dezelfde rest hebben bij de \omterm{thm:g12:arith:euclid}{euclidische
deling} door $n$.
\end{definition}

\begin{proposition}[Verenigbaarheid met de bewerkingen]\label{prop:g12:arith:congops}
Geldt $a \equiv b \pmod n$ en $c \equiv d \pmod n$, dan is
\[
a + c \equiv b + d, \qquad
ac \equiv bd, \qquad
a^k \equiv b^k \ (k \in \N) \pmod n .
\]
\end{proposition}

\begin{proof}
$n$ \omterm{def:g12:arith:divides}{deelt} $(a-b) + (c-d) = (a+c) - (b+d)$, en ook
$ac - bd = a(c - d) + d(a - b)$ is een veelvoud van $n$. De regel voor de
machten volgt met inductie uit die voor het product.
\end{proof}

\begin{method}[Machten modulo $n$ berekenen]\label{met:g12:arith:powers}
Om $a^k \bmod n$ te berekenen, herleid je eerst het grondtal modulo $n$, zoek
je daarna een kleine macht van $a$ die congruent is met $\pm1$, en gebruik je
die om de exponent in te korten. Bijvoorbeeld $2^{100} \bmod 7$: omdat
$2^3 = 8 \equiv 1 \pmod 7$ en $100 = 3\times33 + 1$,
\[
2^{100} = \left(2^{3}\right)^{33} \times 2 \equiv 1^{33}\times 2 = 2 \pmod 7 .
\]
\end{method}

\section{Ggd, Bézout en Gauss}

\begin{definition}[Ggd]\label{def:g12:arith:gcd}
Zij $a$ en $b$ \omterm{def:g10:numbers:sets}{gehele getallen} die niet allebei nul zijn. De \emph{grootste
gemene deler}\index{grootste gemene deler} $\gcd(a, b)$ is het grootste gehele
getal dat zowel $a$ als $b$ \omterm{def:g12:arith:divides}{deelt}. Is $\gcd(a,b) = 1$, dan heten $a$ en $b$
\emph{relatief priem}\index{relatief priem}.
\end{definition}

\begin{proposition}[Algoritme van Euclides]\label{prop:g12:arith:euclidalgo}
\index{algoritme van Euclides}
Geldt $a = bq + r$ ($b \neq 0$), dan is $\gcd(a, b) = \gcd(b, r)$. Door de
\omterm{thm:g12:arith:euclid}{euclidische deling} te herhalen bereken je dus $\gcd(a,b)$: de ggd is de
laatste rest verschillend van nul.
\end{proposition}

\begin{proof}
Elke gemene deler van $a$ en $b$ \omterm{def:g12:arith:divides}{deelt} $r = a - bq$
(\cref{prop:g12:arith:divprops}) en is dus een gemene deler van $b$ en $r$;
en omgekeerd, want $a = bq + r$. De twee paren hebben dezelfde gemene delers,
dus dezelfde ggd. Het algoritme stopt omdat de resten een strikt dalende \omterm{def:g12:seq:sequence}{rij}
van niet-negatieve \omterm{def:g10:numbers:sets}{gehele getallen} vormen.
\end{proof}

\begin{example}\label{ex:g12:arith:euclidalgo}
$\gcd(252, 198)$: $252 = 198 + 54$; $198 = 3\times54 + 36$;
$54 = 36 + 18$; $36 = 2 \times 18 + 0$. Dus $\gcd(252,198) = 18$.
\end{example}

\begin{theorem}[Identiteit van Bézout]\label{thm:g12:arith:bezout}
\index{identiteit van Bézout}
Zij $a$ en $b$ \omterm{def:g10:numbers:sets}{gehele getallen} die niet allebei nul zijn, en $d = \gcd(a,b)$.
Er bestaan $u, v \in \Z$ met
\[
au + bv = d .
\]
In het bijzonder zijn $a$ en $b$ \omterm{def:g12:arith:gcd}{relatief priem} als en slechts als
$au + bv = 1$ voor zekere \omterm{def:g10:numbers:sets}{gehele getallen} $u$ en $v$.
\end{theorem}

\begin{proof}
Doorloop het \omterm{prop:g12:arith:euclidalgo}{algoritme van Euclides} achterwaarts: elke rest is een gehele
combinatie van de twee voorgaande, en de begingegevens $a$ en $b$ zijn
combinaties van zichzelf; door aflopend te substitueren is de laatste rest $d$
verschillend van nul een gehele combinatie van $a$ en $b$. (In
\cref{ex:g12:arith:euclidalgo}:
$18 = 54 - 36 = 54 - (198 - 3\times54) = 4\times54 - 198 =
4(252 - 198) - 198 = 4\times252 - 5\times198$.)

Voor de gelijkwaardigheid: is $\gcd(a,b) = 1$, dan levert Bézout $u$ en $v$;
omgekeerd \omterm{def:g12:arith:divides}{deelt} elke gemene deler van $a$ en $b$ ook $au + bv = 1$, zodat
$\gcd(a,b) = 1$.
\end{proof}

\begin{theorem}[Lemma van Gauss]\label{thm:g12:arith:gauss}
\index{lemma van Gauss}
Zij $a, b, c \in \Z$. Geldt $a \mid bc$ en $\gcd(a, b) = 1$, dan is
$a \mid c$.
\end{theorem}

\begin{proof}
Bézout geeft $au + bv = 1$; vermenigvuldig met $c$: $acu + bcv = c$. Beide
termen van het linkerlid zijn veelvouden van $a$ (de tweede omdat
$a \mid bc$), en dus is $c$ dat ook.
\end{proof}

\begin{corollary}\label{cor:g12:arith:coprimeprod}
Geldt $a \mid c$, $b \mid c$ en $\gcd(a,b) = 1$, dan is $ab \mid c$.
\end{corollary}

\begin{proof}
Schrijf $c = ak$. Uit $b \mid ak$ en $\gcd(a,b)=1$ geeft Gauss $b \mid k$,
zeg $k = bl$; dan is $c = abl$.
\end{proof}

\section{Priemgetallen}

\begin{definition}[Priem]\label{def:g12:arith:prime}
Een \omterm{def:g10:numbers:sets}{geheel getal} $p \geq 2$ heet \emph{priem}\index{priemgetal} als zijn enige
positieve delers $1$ en $p$ zijn.
\end{definition}

\begin{proposition}\label{prop:g12:arith:primedivides}
Elk \omterm{def:g10:numbers:sets}{geheel getal} $n \geq 2$ heeft een priemdeler; is $n$ niet \omterm{def:g12:arith:prime}{priem}, dan heeft
het een priemdeler $\leq \sqrt n$. \omterm{def:g12:arith:divides}{Deelt} een \omterm{def:g12:arith:prime}{priemgetal} $p$ een product $ab$,
dan is $p \mid a$ of $p \mid b$ (\emph{lemma van Euclides}).
\end{proposition}

\begin{proof}
De kleinste deler $d \geq 2$ van $n$ is \omterm{def:g12:arith:prime}{priem} (elke echte deler van $d$ zou
een kleinere deler van $n$ zijn). Is $n = de$ samengesteld met
$2 \leq d \leq e$, dan is $d^2 \leq de = n$, dus $d \leq \sqrt n$. Voor het
lemma van Euclides: geldt $p \nmid a$, dan is $\gcd(p, a) = 1$ (de enige
delers van $p$ zijn $1$ en $p$), en het lemma van Gauss geeft $p \mid b$.
\end{proof}

\begin{theorem}[Euclides]\label{thm:g12:arith:infiniteprimes}
Er zijn oneindig veel priemgetallen.
\end{theorem}

\begin{proof}
Neem een willekeurige eindige lijst $p_1, \dots, p_k$ van priemgetallen en
beschouw $N = p_1 p_2 \cdots p_k + 1$. Een zeker \omterm{def:g12:arith:prime}{priemgetal} $p$ \omterm{def:g12:arith:divides}{deelt} $N$;
maar geen enkele $p_i$ \omterm{def:g12:arith:divides}{deelt} $N$ (de rest is $1$), dus $p$ is een \omterm{def:g12:arith:prime}{priemgetal}
dat niet in de lijst staat. Geen enkele eindige lijst put de priemgetallen
uit.
\end{proof}

\begin{theorem}[Hoofdstelling van de rekenkunde]\label{thm:g12:arith:fta}
\index{hoofdstelling van de rekenkunde}
Elk \omterm{def:g10:numbers:sets}{geheel getal} $n \geq 2$ is een product van priemgetallen, en die
ontbinding is uniek op de volgorde van de factoren na:
\[
n = p_1^{\alpha_1} p_2^{\alpha_2} \cdots p_r^{\alpha_r},
\qquad p_1 < p_2 < \dots < p_r \text{ priem},\ \alpha_i \geq 1 .
\]
\end{theorem}

\begin{proof}
\emph{Bestaan}, met sterke inductie: is $n$ \omterm{def:g12:arith:prime}{priem}, dan is het zijn eigen
ontbinding; anders is $n = de$ met $2 \leq d, e < n$, en beide \omterm{def:g10:algebra:expand}{ontbinden}
volgens de inductiehypothese. \emph{Uniciteit}: onderstel
$p_1\cdots p_s = q_1 \cdots q_t$ (priemgetallen, herhalingen toegestaan).
Volgens het lemma van Euclides \omterm{def:g12:arith:divides}{deelt} $p_1$ een zekere $q_j$, en omdat $p_1$
\omterm{def:g12:arith:prime}{priem} is, is $p_1 = q_j$; schrap en herhaal. De twee ontbindingen stemmen term
voor term overeen.
\end{proof}

\begin{theorem}[Kleine stelling van Fermat]\label{thm:g12:arith:fermat}
\index{kleine stelling van Fermat}
Zij $p$ \omterm{def:g12:arith:prime}{priem} en $a \in \Z$ met $p \nmid a$. Dan is
\[
a^{p-1} \equiv 1 \pmod p .
\]
Voor elke $a \in \Z$ (zonder enige voorwaarde) geldt $a^p \equiv a \pmod p$.
\end{theorem}

\begin{proof}
Beschouw de $p - 1$ \omterm{def:g10:numbers:sets}{gehele getallen} $a, 2a, 3a, \dots, (p-1)a$ modulo $p$.
Geen ervan is $\equiv 0$ (geldt $p \mid ka$ met $1 \leq k \leq p-1$, dan
dwingt het lemma van Euclides $p \mid k$ af, wat onmogelijk is), en ze zijn
paarsgewijs verschillend modulo $p$ (geldt $ka \equiv la$, dan is
$p \mid (k - l)a$, dus $p \mid k - l$, dus $k = l$). Modulo $p$ zijn het dus
de getallen $1, 2, \dots, p-1$ in een zekere volgorde. Vermenigvuldig alle
\omterm{def:g12:arith:congruence}{congruenties}:
\[
a^{p-1}\,(p-1)! \equiv (p-1)! \pmod p .
\]
Omdat $p$ geen enkele van $1, \dots, p-1$ \omterm{def:g12:arith:divides}{deelt}, mag je met herhaald gebruik
van het lemma van Euclides $(p-1)!$ wegdelen, en blijft $a^{p-1} \equiv 1$
over. De tweede vorm volgt door met $a$ te vermenigvuldigen (en is triviaal
zodra $p \mid a$).
\end{proof}

\begin{example}[Toepassing in de cryptografie]\label{ex:g12:arith:rsa}
De stelling van Fermat maakt het machtsverheffen modulo $n$ omkeerbaar zodra
de exponenten geschikt gekozen zijn --- het hart van het
\emph{RSA}-cryptosysteem. Met grote priemgetallen $p$ en $q$ en $n = pq$
publiceer je $n$ en een exponent $e$; versleutelen is $x \mapsto x^e \bmod n$.
Ontsleutelen vraagt een exponent $d$ met $ed \equiv 1 \pmod{(p-1)(q-1)}$, en
die kan alleen berekend worden door wie $p$ en $q$ kent --- en $p$ en $q$ uit
$n$ terugvinden betekent een getal van honderden cijfers \omterm{def:g10:algebra:expand}{ontbinden}, wat geen
enkel bekend algoritme in redelijke tijd doet.
\end{example}

\section{Oefeningen}

\begin{exercise}[$\star$]\label{exo:g12:arith:1}
Bereken het quotiënt en de rest van de \omterm{thm:g12:arith:euclid}{euclidische deling} van $2026$ door
$17$, en van $-2026$ door $17$.
\end{exercise}

\begin{exercise}[$\star$]\label{exo:g12:arith:2}
Wat is de rest van $7^{100}$ modulo $10$? (Wat is het laatste cijfer van
$7^{100}$?)
\end{exercise}

\begin{exercise}[$\star$]\label{exo:g12:arith:3}
Bereken met het \omterm{prop:g12:arith:euclidalgo}{algoritme van Euclides} $\gcd(1071, 462)$, en zoek \omterm{def:g10:numbers:sets}{gehele
getallen} $u$ en $v$ met $1071u + 462v = \gcd(1071, 462)$.
\end{exercise}

\begin{exercise}[$\star$]\label{exo:g12:arith:4}
Toon aan dat $n^2$ voor elke $n \in \Z$ congruent is met $0$ of $1$ modulo
$4$. Leid af dat een \omterm{def:g10:numbers:sets}{geheel getal} $\equiv 3 \pmod 4$ nooit een som van twee
kwadraten is.
\end{exercise}

\begin{exercise}[$\star\star$]\label{exo:g12:arith:5}
Toon aan dat $n(n+1)(2n+1)$ voor alle $n \in \N$ deelbaar is door $6$.
\end{exercise}

\begin{exercise}[$\star\star$]\label{exo:g12:arith:6}
Los in $\Z$ de \omterm{def:g12:arith:congruence}{congruentie} $5x \equiv 3 \pmod{11}$ op.
(Tip: zoek de inverse van $5$ modulo $11$.)
\end{exercise}

\begin{exercise}[$\star\star$]\label{exo:g12:arith:7}
Los in $\Z \times \Z$ de diofantische \omterm{def:g10:algebra:equation}{vergelijking}
\[
17x - 40y = 1
\]
op, en beschrijf daarna alle oplossingen van $17x - 40y = 6$.
\end{exercise}

\begin{exercise}[$\star\star$]\label{exo:g12:arith:8}
Toon aan dat $\sqrt2$ \omterm{ex:g10:numbers:classify}{irrationaal} is met behulp van de uniciteit van de
priemontbinding (vergelijk de exponent van $2$ in beide leden van
$a^2 = 2b^2$).
\end{exercise}

\begin{exercise}[$\star\star\star$]\label{exo:g12:arith:9}
Zij $p$ een \omterm{def:g12:arith:prime}{priemgetal}.
\begin{enumerate}
  \item Toon aan dat $p$ voor $1 \leq k \leq p - 1$ het getal $\dbinom{p}{k}$
        \omterm{def:g12:arith:divides}{deelt}. (Tip: gebruik $k\binom pk = p\binom{p-1}{k-1}$,
        \cref{exo:g12:comb:7}, en het lemma van Gauss.)
  \item Leid daaruit met inductie op $a \geq 0$ een tweede bewijs af van de
        kleine stelling van Fermat in de vorm $a^p \equiv a \pmod p$.
\end{enumerate}
\end{exercise}

\begin{exercise}[$\star\star\star$]\label{exo:g12:arith:10}
\emph{(Chinees restprobleem.)} Zoek alle \omterm{def:g10:numbers:sets}{gehele getallen} $n$ waarvoor
\[
n \equiv 2 \pmod 3, \qquad n \equiv 3 \pmod 5, \qquad n \equiv 2 \pmod 7 .
\]
(Tip: los eerst de twee eerste voorwaarden op en voeg daarna de derde toe;
coëfficiënten van Bézout helpen.)
\end{exercise}

\section{Opgave: geheime codes en controlecijfers}

\begin{problem}[{Weekendopgave --- congruenties bewaken elke streepjescode en
elke bankkaart, en de kleine stelling van Fermat bedient het slot op de
geheimen van de wereld}]
\label{pb:g12:arith:1}
G.\,H.~Hardy schepte er in 1940 over op dat de getaltheorie ``onbezoedeld''
door toepassingen was. Tachtig jaar later spreekt elke piep aan de kassa, elke
kaartbetaling en elk versleuteld bericht hem tegen --- en dat met precies het
gereedschap van dit hoofdstuk: \omterm{def:g12:arith:congruence}{congruenties}
(\cref{prop:g12:arith:congops}), inversen van Bézout
(\cref{thm:g12:arith:bezout}) en de kleine stelling van Fermat
(\cref{exo:g12:arith:9}). Deze opgave controleert de codes, breekt een
speelgoedversie van het slot open, en leert waarom het echte slot standhoudt.

\medskip
\noindent\textbf{Deel I --- Vlot met \omterm{def:g12:arith:congruence}{congruenties}.}
\begin{enumerate}
  \item Bereken $2026 \bmod 7$; en daarna het laatste cijfer van $7^{100}$
        (zoek de cyclus van de machten van $7$ modulo $10$).
  \item Snel machtsverheffen (\cref{met:g12:arith:powers}): bereken
        $5^{117} \bmod 13$ (vertrek van $5^2 \equiv -1$).
  \item Los $3x \equiv 5 \pmod 7$ op.
  \item Pas het \omterm{prop:g12:arith:euclidalgo}{algoritme van Euclides} toe op $(97, 35)$, substitueer
        achterwaarts om \omterm{def:g10:numbers:sets}{gehele getallen} $u$ en $v$ met $97u + 35v = 1$ te
        vinden, en leid de inverse van $35$ modulo $97$ af.
  \item Formuleer precies wanneer $a$ inverteerbaar is modulo $n$, en welke
        stelling de inverse levert.
\end{enumerate}

\medskip
\noindent\textbf{Deel II --- Controlecijfers.}
\begin{enumerate}[resume]
  \item ISBN-10: de tien cijfers $d_1 \dots d_{10}$ van een boekcode moeten
        voldoen aan
        $10d_1 + 9d_2 + \dots + 2d_9 + 1d_{10} \equiv 0 \pmod{11}$. Ga het
        echte ISBN $0\,306\,40615\,2$ na.
  \item Bewijs dat het ISBN-schema \emph{elke} fout in één cijfer opspoort:
        verandert één cijfer met $d \not\equiv 0$, dan verandert de gewogen
        som met $w d$ met $1 \leq w \leq 10$ --- waarom kan dat nooit
        $\equiv 0 \pmod{11}$ zijn (\cref{thm:g12:arith:gauss})?
  \item Bewijs dat het ook elke verwisseling van twee naburige (verschillende)
        cijfers opspoort. Leg daarna het geheim van het ontwerp uit: welke
        eigenschap van $11$ deed beide bewijzen werken, en wat kan er misgaan
        met \omterm{def:g12:complex:modulus}{modulus} $10$?
  \item Streepjescodes EAN-13 wegen de cijfers met $1, 3, 1, 3, \dots$ modulo
        $10$. Bereken het controlecijfer dat $978\,2940199\,05$ vervolledigt.
        Welke verwisselingen van naburige cijfers spoort EAN \emph{niet} op?
        (Wanneer is $2(a - b) \equiv 0 \pmod{10}$?)
  \item Bankkaarten gebruiken het schema van Luhn: verdubbel van rechts af
        elk tweede cijfer (en trek $9$ af zodra het dubbele boven $9$ uitkomt),
        tel alles op, en eis een veelvoud van $10$. Ga het testnummer
        $4539\,1488\,0343\,6467$ na.
  \item In één zin: wat leverde de priemmodulus het ISBN op dat EAN en Luhn,
        geketend aan $10$, niet kunnen hebben?
\end{enumerate}

\medskip
\noindent\textbf{Deel III --- Het slot van Fermat.}
\begin{enumerate}[resume]
  \item Een valstrik vóór de schat: bereken $2^{10} \bmod 341$, leid
        $2^{340} \bmod 341$ af --- en ontbind daarna $341$. Wat zegt dit
        voorbeeld (een \emph{pseudopriemgetal van Fermat}) over het gebruik van
        de kleine stelling van Fermat als priemtest?
  \item RSA in het klein: neem $p = 3$ en $q = 11$, dus $n = 33$ en
        $(p-1)(q-1) = 20$; de publieke exponent is $e = 3$. Zoek de private
        exponent $d$ met $3d \equiv 1 \pmod{20}$ (de methode van vraag 4).
  \item Versleutel het bericht $m = 4$: bereken $c = m^3 \bmod 33$.
  \item Ontsleutel: bereken $c^d \bmod 33$ (gebruik
        $c \equiv -2 \pmod{33}$) en haal het bericht terug.
  \item Waarom het ontsleutelen altijd werkt: toon aan dat
        $m^{21} \equiv m$ zowel modulo $3$ als modulo $11$ (de kleine stelling
        van Fermat in elk van beide werelden), en besluit modulo $33$
        (\cref{thm:g12:arith:gauss} lijmt de twee \omterm{def:g12:arith:congruence}{congruenties} aan elkaar).
        Waar kwam de bijzondere vorm $1 + 20k$ van $21 = ed$ binnen?
  \item De veiligheid van het slot: iedereen kent $n$ en $e$; om $d$ terug te
        vinden heb je $(p-1)(q-1)$ nodig, en dus de factoren van $n$. Onze $33$
        ontbindt in één oogopslag --- waarom beschermt hetzelfde schema, met
        een $n$ van zeshonderd cijfers, de banken van de wereld? (Eén zin over
        de asymmetrie tussen vermenigvuldigen en \omterm{def:g10:algebra:expand}{ontbinden}.)
\end{enumerate}

\medskip
\noindent\textbf{Deel IV --- Klassiekers.}
\begin{enumerate}[resume]
  \item De oude Chinese soldatentelling (vergelijk met
        \cref{exo:g12:arith:10}): een aantal soldaten laat rest $2$ bij
        opstelling per $3$ en rest $3$ bij opstelling per $5$. Zoek alle
        mogelijke aantallen, en leg uit waarom het antwoord uniek is modulo
        $15$.
  \item Eindelijk bewijzen van één regel: leid uit $10 \equiv 1 \pmod 9$ af
        dat elk getal congruent is met zijn cijfersom modulo $9$; leid uit
        $10 \equiv -1 \pmod{11}$ de regel van de alternerende som voor $11$
        af. (Het onderbouwvolume bewees die met expliciete algebra --- bewonder
        de beknoptheid.)
  \item Slotstuk --- Hardy tegen de streepjescode: vat het gereedschap van het
        hoofdstuk samen (rekenen met \omterm{def:g12:arith:congruence}{congruenties}, inversen van Bézout, de
        kleine stelling van Fermat, het lijmen van relatief prieme moduli) en
        zeg waar elk stuk in deze opgave op zijn plaats klikte; geef daarna het
        moderne oordeel over ``onbezoedeld''.
\end{enumerate}
\end{problem}
